 
\documentclass[12pt]{article}



\usepackage{amssymb}
\usepackage{amsmath}
\usepackage{ifthen}
\usepackage[table]{xcolor}
\usepackage{minitoc}
\usepackage{array}
\usepackage{graphicx}

\definecolor{yellow}{cmyk}{0,0,1,0}
\renewcommand{\arraystretch}{1.4}
\newcommand{\R}{\mathbb{R}}

\usepackage{colortbl}
\usepackage[top=0.75in, bottom=0.75in, left=1.25in, right=1.25in]{geometry}

%% Page size
%\setlength{\oddsidemargin}{13pt}
%%\setlength{\evensidemargin}{-0.5in}
%\setlength{\textheight}{674pt}
%\setlength{\textwidth}{426pt}
%\setlength{\topmargin}{-10pt}
%\setlength{\voffset}{-30pt}

\renewcommand{\arraycolsep}{3pt}


\input{color_flatex}


% Page size
%\setlength{\oddsidemargin}{-0.25in}
%\setlength{\evensidemargin}{-0.25in}
%\setlength{\textheight}{9.25in}
%\setlength{\textwidth}{7in}
%\setlength{\topmargin}{-0.5in}


% ----------------------------------------------------------------
\title{HPC Final \\ CS 378 Programming for Correctness and Performance}

\author{Umer Khan, Ahmad Raeq Ahsen, Saood Usmani}

\begin{document}
	\maketitle

\section{Operation}

Consider the operation
\[
B := L^T B 
\]
where $ L^T $ is a $ m \times m $ lower triangular matrix that will be transposed and behave like an upper triangle and $ B $ is a $ m \times n $ matrix.
This is a special case of  
matrix-matrix multiplication, 
with the {\sc l}ower triangular matrix on the {\sc l}eft, 
We will refer to this operation
as {\sc Trmm\_lltn} where the {\sc lltn} stands for
\underline{l}eft
\underline{l}ower
\underline{t}ranspose
\underline{n}onunit diagonal.
The {\sc n}onunit diagonal means we will use the entries of the matrix that are stored on the diagonal.  In some other cases, the entries on the diagonal may implicitly be taked to all equal one ({\sc u}nit).

\section{Precondition and Postcondition}

In the precondition 
\[
B = \widehat B
\]
$ \widehat B $ denotes the original contents of $ B $.
This allows us to express the state upon completion, the postcondition, as
\[
B = L^T \widehat B.
\]
It is implicitly assumed that $ L $ is nonunit lower triangular.

\section{Purpose (BONUS)}
This operation can be used when solving a set of linear system of equations with a symmetric matrix stored in the lower triangle ($A = LL^T$ or $A - LDL^T$). This is because the equation needs both Solve Ly=b via forward substitution and $L^Tx = y$ via backward substitution.
\section{Partitioned Matrix Expressions and Loop Invariants}

There are two PMEs for this operation.

\subsection{PME 1}

To derive the first PME, partition
\[
L \rightarrow
\left(\begin{array}{c I c}
	L_{TL} & L_{BL}^T \\ \whline
	0 & L_{BR}
\end{array}\right)
,
\quad \mbox{and} \quad
B \rightarrow 
\left(\begin{array}{c}
	B_T \\ \whline
	B_B
\end{array}\right).
\]
Substituting these into the postcondition
yields
\[
\left(\begin{array}{c}
	B_T \\ \whline
	B_B
\end{array}\right)
=
\left(\begin{array}{c I c}
	L_{TL} & L_{BL}^T \\ \whline
	0 & L_{BR}
\end{array}\right)
\left(\begin{array}{c}
	\widehat B_T \\ \whline
	\widehat B_B
\end{array}\right)
\]
or, equivalently,
\[
\left(\begin{array}{c}
	B_T \\ \whline
	B_B
\end{array}\right)
=
\left(\begin{array}{c}
	L_{TL} \widehat B_T + L_{BL}^T \widehat B_B \\ \whline
	L_{BR} \widehat B_B
\end{array}\right),
\]
which we refer to as the first PME for this operations.

From this, we can choose two loop invariants:
\begin{description}
	\item
	{\bf Invariant 1:}
\begin{equation}
	\left(\begin{array}{c}
		B_T \\ \whline
		B_B
	\end{array}\right)
	= 
	\left(\begin{array}{c}
		L_{TL}\widehat B_T \\ \whline
		\widehat B_B
	\end{array}\right).	
\label{eq:inv1}
\end{equation}	
 \\
	(The bottom part has been left alone and the top part has been partially computed).
	\item
	{\bf Invariant 2:}
\begin{equation}
	\left(\begin{array}{c}
		B_T \\ \whline
		B_B
	\end{array}\right) = 
	\left(\begin{array}{c}
		L_{TL} \widehat B_T + (L_{BL}^T)^T \widehat B_B  \\ \whline
		\widehat B_B
	\end{array}\right).
\label{eq:inv2}
	\end{equation} \\
	(The bottom part has been left alone and the top part has been completely computed).
\end{description}
% TODO: 
\subsection{PME 2}

To derive the second PME, partition
\[
B \rightarrow \left(\begin{array}{c I c}
	B_L & B_R 
\end{array}\right)
\]
and do not partition $ L^T $.
Substituting these into the postcondition
yields
\[
\left(\begin{array}{c I c}
	B_L & B_R 
\end{array}\right)
=
L^T 
\left(\begin{array}{c I c}
	\widehat B_L & \widehat B_R 
\end{array}\right)
\]
or, equivalently,
\[
\left(\begin{array}{c I c}
	B_L & B_R 
\end{array}\right)
=
\left(\begin{array}{c I c}
	L^T \widehat B_L & L^T \widehat B_R 
\end{array}\right),
\]
which we refer to as the second PME.

From this, we can choose two more loop invariants:
\begin{description}
	\item
	{\bf Invariant 3:}

\begin{equation}
		\left(\begin{array}{c I c}
		B_L & B_R 
	\end{array}\right) =
	\left(\begin{array}{c I c}
		L^T \widehat B_L & \widehat B_R 
	\end{array}\right).
\label{eq:inv3}
\end{equation}
\\
	(The left part has been completely finished and the right part has been left untouched).
	\item
	{\bf Invariant 4:}
	
\begin{equation}
		\left(\begin{array}{c I c}
		B_L & B_R 
	\end{array}\right) =
	\left(\begin{array}{c I c}
		\widehat B_L & L^T \widehat B_R 
	\end{array}\right).
\label{eq:inv4}
\end{equation}\\
	(The right part has been completely finished and the left part has been left untouched).
\end{description}


\section{Deriving the Algorithms}

\subsection{Loop Invariant 3}

The unblocked algorithm for loop invariant 3 is showing in Figure~\ref{fig:unb_inv1}. While the blocked algorithm for loop invariant 3 is showing in Figure~\ref{fig:blk_inv3}.


\resetsteps


\resetsteps      % Reset all the commands to create a blank worksheet  

% Define the operation to be computed

\renewcommand{\operation}{ \left[ B \right] := \mbox{\sc Trmm\_lltn\_unb\_var1}( L, B ) }

\renewcommand{\routinename}{\operation}

% Step 1a: Precondition 

\renewcommand{\precondition}{
  B = \widehat{B}
}

% Step 1b: Postcondition 

\renewcommand{\postcondition}{ 
  \left[B \right]
  =
  \mbox{Trmm\_lltn}( L, \widehat{B} )
}

% Step 2: Invariant 
% Note: Right-hand side of equalities must be updated appropriately

\renewcommand{\invariant}{
  \left(\begin{array}{c I c}
     B_L & B_R
  \end{array}\right)  = 
  \left(\begin{array}{c I c}
     L^T \widehat B_L   & \widehat B_R
  \end{array}\right)
}

% Step 3: Loop-guard 

\renewcommand{\guard}{
  n( B_L ) < n( B )
}

% Step 4: Initialize 

\renewcommand{\partitionings}{
  $
  B \rightarrow
  \left(\begin{array}{c I c}
     B_L & B_R
  \end{array}\right)
  $
}

\renewcommand{\partitionsizes}{
  $ B_L $ has $ 0 $ columns
}

% Step 5a: Repartition the operands 

\renewcommand{\repartitionings}{
  $  \left(\begin{array}{c I c}
     B_L & B_R
  \end{array}\right)
  \rightarrow
  \left(\begin{array}{c I c c}
     B_0 & b_1 & B_2
  \end{array}\right)
  $
}

\renewcommand{\repartitionsizes}{
  $ b_1 $ has $ 1 $ column}

% Step 5b: Move the double lines 

\renewcommand{\moveboundaries}{
$  
  B \rightarrow
  \left(\begin{array}{c I c c}
     B_L & B_R
  \end{array}\right)
  \leftarrow
  \left(\begin{array}{c c I c}
     B_0 & b_1 & B_2
  \end{array}\right)
  $
}

% Step 6: State before update
% Note: The below needs editing consistent with loop-invariant!!!

\renewcommand{\beforeupdate}{$B = \left(\begin{array}{c I c c}
     L \widehat B_0 & b_1  B_2
  \end{array}\right)
   $}


% Step 7: State after update
% Note: The below needs editing consistent with loop-invariant!!!

\renewcommand{\afterupdate}{$B=\left(\begin{array}{c c I c}
     L \widehat B_0 & L \widehat b_1 & B_2
  \end{array}\right)$}


% Step 8: Insert the updates required to change the 
%         state from that given in Step 6 to that given in Step 7
% Note: The below needs editing!!!

\renewcommand{\update}{
$
  \begin{array}{l}          % do not delete this line 
    
    \mbox{$b_1 := L b_1$} \\ % replace \mbox{...} with update line but leave 
  \end{array}               % do not delete this line 
$
}




\begin{figure}
	\begin{center}
		\FlaWorksheet
	\end{center}
	\caption{Unblocked Algorithm for Loop Invariant 3}
	\label{fig:unb_inv1}
\end{figure}


\resetsteps

\input{derivations/Trmm_lltn_blk_var1}

\begin{figure}
	\begin{center}
		\FlaWorksheet
	\end{center}
	\caption{Blocked Algorithm for Loop Invariant 3}
	\label{fig:blk_inv3}
\end{figure}


\newpage


\section{Operation}

Consider the operation  
\[
C := A A^T + C
\]
where \( A \) is a \( m \times k \) matrix and \( C \) is a \( m \times m \) symmetric matrix, stored in its lower-triangular part. This is a special case of  
\textit{symmetric rank-\( k \)} update,  
with the matrix \( A \) on the \textsc{l}eft and its transpose \( A^T \) on the right (And by definition of symmetric, we know that $AA^T$ is symmetric).  
We refer to this operation as  
\textsc{Syrk\_ln}, where the \textsc{ln} stands for:
\begin{itemize}
  \item \underline{l}eft (the matrix \( A \) is on the left)
  \item \underline{n}o transpose (the matrix \( A \) is used as-is, without transposition)
\end{itemize}

This operation updates the symmetric matrix \( C \) using the product \( A A^T \), storing the result only in the lower triangular part to exploit symmetry and save space.

\section{Precondition and Postcondition}

In the precondition 
\[
C = \widehat C
\]
$ \widehat C $ denotes the original contents of $ C $.
This allows us to express the state upon completion, the postcondition, as
\[
C = A A^T + \widehat{C}
\]
\section{Purpose (BONUS)}
In computing Gram matrices, the operation efficiently accumulates inner products (which is useful for Gram Matrices) between data vectors, exploiting symmetry to reduce computation and storage. It is commonly used in kernel methods and PCA for similarity or covariance estimation.
\section{Partitioned Matrix Expressions and Loop Invariants}

\subsection{PME 1}

To derive the first PME, partition
\[
A \rightarrow
\left(\begin{array}{c}
	A_T \\ \whline
	A_B
\end{array}\right),
\quad \text{and} \quad
\widehat{C} \rightarrow
\left(\begin{array}{c I c}
	\widehat{C}_{TL} & \widehat{C}_{BL}^T \\ \whline
	\widehat{C}_{BL} & \widehat{C}_{BR}
\end{array}\right).
\]

Substituting into the postcondition yields
\[
C = AA^T + \widehat{C}
\Rightarrow
\left(\begin{array}{c I c}
	C_{TL} & C_{BL}^T \\ \whline
	C_{BL} & C_{BR}
\end{array}\right)
=
\left(\begin{array}{c}
	A_T \\ \whline
	A_B
\end{array}\right)
\left(\begin{array}{c}
	A_T \\ \whline
	A_B
\end{array}\right)^T
+
\left(\begin{array}{c I c}
	\widehat{C}_{TL} & \widehat{C}_{BL}^T \\ \whline
	\widehat{C}_{BL} & \widehat{C}_{BR}
\end{array}\right)
\]

Evaluating the multiplication yields:
\[
\left(\begin{array}{c I c}
	C_{TL} & C_{BL}^T \\ \whline
	C_{BL} & C_{BR}
\end{array}\right)
=
\left(\begin{array}{c I c}
	A_T A_T^T + \widehat{C}_{TL} & A_T A_B^T + \widehat{C}_{BL}^T \\ \whline
	A_B A_T^T + \widehat{C}_{BL} & A_B A_B^T + \widehat{C}_{BR}
\end{array}\right)
\]

This is the first PME for the operation.

\begin{description}
  \item
  \textbf{Invariant 1:}
  \begin{equation}
    \left(\begin{array}{c I c}
      C_{TL} & C_{BL}^T \\ \whline
      C_{BL} & C_{BR}
    \end{array}\right)
    =
    \left(\begin{array}{c I c}
      A_T A_T^T + \widehat{C}_{TL} & \widehat{C}_{BL}^T \\ \whline
      \widehat{C}_{BL} & \widehat{C}_{BR}
    \end{array}\right)
  \label{eq:inv1}
  \end{equation}
  \\
  (Only the top-left block has been  updated.)

  \item
  \textbf{Invariant 2:}
  \begin{equation}
    \left(\begin{array}{c I c}
      C_{TL} & C_{BL}^T \\ \whline
      C_{BL} & C_{BR}
    \end{array}\right)
    =
    \left(\begin{array}{c I c}
      A_T A_T^T + \widehat{C}_{TL} & A_T A_B^T + \widehat{C}_{BL}^T \\ \whline
      A_B A_T^T + \widehat{C}_{BL} & \widehat{C}_{BR}
    \end{array}\right)
  \label{eq:inv2}
  \end{equation}
  \\
  (The entire top row and left column have been fully updated.)

  \item
  \textbf{Invariant 3:}
  \begin{equation}
    \left(\begin{array}{c I c}
      C_{TL} & C_{BL}^T \\ \whline
      C_{BL} & C_{BR}
    \end{array}\right)
    =
    \left(\begin{array}{c I c}
      \widehat{C}_{TL} & \widehat{C}_{BL}^T \\ \whline
      \widehat{C}_{BL} & A_B A_B^T + \widehat{C}_{BR}
    \end{array}\right)
  \label{eq:inv3}
  \end{equation}
  \\
  (Only the bottom-right block has been updated.)

  \item
  \textbf{Invariant 4:}
  \begin{equation}
    \left(\begin{array}{c I c}
      C_{TL} & C_{BL}^T \\ \whline
      C_{BL} & C_{BR}
    \end{array}\right)
    =
    \left(\begin{array}{c I c}
      \widehat{C}_{TL} & A_T A_B^T + \widehat{C}_{BL}^T \\ \whline
      A_B A_T^T + \widehat{C}_{BL} & A_B A_B^T + \widehat{C}_{BR}
    \end{array}\right)
  \label{eq:inv4}
  \end{equation}
  \\
  (The entire bottom row and right column have been updated, leaving the top-left untouched.)
\end{description}

\subsection{PME 2}

To derive the second PME, partition
\[
A \rightarrow \left(\begin{array}{c I c}
	A_L & A_R 
\end{array}\right)
\]
and do not partition \( \widehat{C} \).

Substituting this into the postcondition yields
\[
C = A A^T + \widehat{C}
\Rightarrow
C =
\left(\begin{array}{c I c}
	A_L & A_R 
\end{array}\right)
\left(\begin{array}{c}
	A_L^T \\ \whline
	A_R^T 
\end{array}\right)
+
\widehat{C}
\]

Evaluating the matrix product gives:
\[
C = A_L A_L^T + A_R A_R^T + \widehat{C}
\]

This expression is referred to as the second PME for this operation.


From this, we can choose two more loop invariants:
\begin{description}
	\item
	{\bf Invariant 5:}

\begin{equation}
		C = A_L A_L^T + \widehat{C}
\label{eq:inv5}
\end{equation}
\\
	(All of C is being partially computed on A moves from left to right).
\item
{\bf Invariant 6:}
	
\begin{equation}
		C = A_R A_R^T + \widehat{C}
\label{eq:inv6}
\end{equation}\\
	(All of C is being partially computed on A moves from right to left).
\end{description}


\section{Deriving the Algorithms}

\subsection{Loop Invariant \eqref{eq:inv1}}

The unblocked algorithm for loop invariant \eqref{eq:inv1} is showing in Figure 3. While the blocked algorithm for loop invariant \eqref{eq:inv1} is showing in Figure~\ref{fig:blk_inv1}.










\resetsteps      % Reset all the commands to create a blank worksheet  

% Define the operation to be computed

\renewcommand{\operation}{ \left[ C \right] := \mbox{\sc Syrk\_ln\_unb\_var1}( A, C ) }

\renewcommand{\routinename}{\operation}

% Step 1a: Precondition 

\renewcommand{\precondition}{
  C = \widehat{C}
}

% Step 1b: Postcondition 

\renewcommand{\postcondition}{ 
  \left[C \right]
  =
  \mbox{Syrk\_ln}( A, \widehat{C} )
}

% Step 2: Invariant 
% Note: Right-hand side of equalities must be updated appropriately

\renewcommand{\invariant}{
C = A_L A_L^T + \widehat{C}
}

% Step 3: Loop-guard 

\renewcommand{\guard}{
  n( A_L ) < n( A )
}

% Step 4: Initialize 

\renewcommand{\partitionings}{
  $
  A \rightarrow
  \left(\begin{array}{c I c}
     A_L & A_R
  \end{array}\right)
  $
,
  $
  A^T \rightarrow
  \left(\begin{array}{c I c}
     A_L^T & A_R^T
  \end{array}\right)
  $
}

\renewcommand{\partitionsizes}{
  $ A_L $ has $ 0 $ columns,
  $ A_L^T $ has $ 0 $ columns
}

% Step 5a: Repartition the operands 

\renewcommand{\repartitionings}{
  $  \left(\begin{array}{c I c}
     A_L & A_R
  \end{array}\right)
  \rightarrow
  \left(\begin{array}{c I c c}
     A_0 & a_1 & A_2
  \end{array}\right)
  $
,
  $  \left(\begin{array}{c I c}
     A_L^T & A_R^T
  \end{array}\right)
  \rightarrow
  \left(\begin{array}{c I c c}
     A_0^T & a_1^T & A_2^T
  \end{array}\right)
  $
}

\renewcommand{\repartitionsizes}{
  $ a_1 $ has $ 1 $ column,
  $ a_1^T $ has $ 1 $ column}

% Step 5b: Move the double lines 

\renewcommand{\moveboundaries}{
$  
  A \rightarrow
  \left(\begin{array}{c I c}
     A_L & A_R
  \end{array}\right)
  \leftarrow
  \left(\begin{array}{c c I c}
     A_0 & a_1 & A_2
  \end{array}\right)
  $
,
$  
  A^T \rightarrow
  \left(\begin{array}{c I c}
     A_L^T & A_R^T
  \end{array}\right)
  \leftarrow
  \left(\begin{array}{c c I c}
     A_0^T & a_1^T & A_2^T
  \end{array}\right)
  $
}

% Step 6: State before update
% Note: The below needs editing consistent with loop-invariant!!!

\renewcommand{\beforeupdate}{$ C = A_0 A_0^T + \widehat{C}$}


% Step 7: State after update
% Note: The below needs editing consistent with loop-invariant!!!

\renewcommand{\afterupdate}{$ C =  \left(\begin{array}{c c}
     A_L & a_1
  \end{array}\right) \left(\begin{array}{c c}
     A_L^T & a_1^T
  \end{array}\right) + \widehat{C} $}


% Step 8: Insert the updates required to change the 
%         state from that given in Step 6 to that given in Step 7
% Note: The below needs editing!!!

\renewcommand{\update}{
$
  \begin{array}{l}          % do not delete this line 
    \mbox{$C := a_1(a_1^T)^T + C$} \\ % replace \mbox{...} with update line but leave 
  \end{array}               % do not delete this line 
$
}




\begin{figure}
	\begin{center}
		\FlaWorksheet
	\end{center}
	\caption{Blocked Algorithm for Loop Invariant \eqref{eq:inv1}}
	\label{fig:blk_inv1}
\end{figure}


\resetsteps      % Reset all the commands to create a blank worksheet  

% Define the operation to be computed

\renewcommand{\operation}{ \left[ C \right] := \mbox{\sc Syrk\_ln\_blk\_var1}( A, A, C ) }

\renewcommand{\routinename}{\operation}

% Step 1a: Precondition 

\renewcommand{\precondition}{
  C = \widehat{C}
}

% Step 1b: Postcondition 

\renewcommand{\postcondition}{ 
  \left[C \right]
  =
  \mbox{Syrk\_ln}( A, A, \widehat{C} )
}

% Step 2: Invariant 
% Note: Right-hand side of equalities must be updated appropriately

\renewcommand{\invariant}{
C = A_L A_L^T + \widehat{C}
}

% Step 3: Loop-guard 

\renewcommand{\guard}{
  n( A_L ) < n( A )
}

% Step 4: Initialize 

\renewcommand{\partitionings}{
  $
  A \rightarrow
  \left(\begin{array}{c I c}
     A_L & A_R
  \end{array}\right)
  $
,
  $
  A \rightarrow
  \left(\begin{array}{c I c}
     A_L & A_R
  \end{array}\right)
  $
}

\renewcommand{\partitionsizes}{
  $ A_L $ has $ 0 $ columns,
  $ A_L $ has $ 0 $ columns
}

% Step 5a: Repartition the operands 

\renewcommand{\repartitionings}{
  $  \left(\begin{array}{c I c}
     A_L & A_R
  \end{array}\right)
  \rightarrow
  \left(\begin{array}{c I c c}
     A_0 & A_1 & A_2
  \end{array}\right)
  $
,
  $  \left(\begin{array}{c I c}
     A_L^T & A_R^T
  \end{array}\right)
  \rightarrow
  \left(\begin{array}{c I c c}
     A_0^T & A_1^T & A_2^T
  \end{array}\right)
  $
}

\renewcommand{\repartitionsizes}{
  $ A_1 $ has $ b $ columns,
  $ A_1^T $ has $ b $ columns}

% Step 5b: Move the double lines 

\renewcommand{\moveboundaries}{
$  
  A \rightarrow
  \left(\begin{array}{c I c}
     A_L & A_R
  \end{array}\right)
  \leftarrow
  \left(\begin{array}{c c I c}
     A_0 & A_1 & A_2
  \end{array}\right)
  $
,
$  
  A \rightarrow
  \left(\begin{array}{c I c}
     A_L^T & A_R^T
  \end{array}\right)
  \leftarrow
  \left(\begin{array}{c c I c}
     A_0^T & A_1^T & A_2^T
  \end{array}\right)
  $
}

% Step 6: State before update
% Note: The below needs editing consistent with loop-invariant!!!

\renewcommand{\beforeupdate}{$ C = A_0 A_0^T + \widehat{C}$}


% Step 7: State after update
% Note: The below needs editing consistent with loop-invariant!!!

\renewcommand{\afterupdate}{$ C =  \left(\begin{array}{c I c}
     A_L & A_1
  \end{array}\right) \left(\begin{array}{c I c}
     A_L^T & A_1^T
  \end{array}\right) + \widehat{C} $}


% Step 8: Insert the updates required to change the 
%         state from that given in Step 6 to that given in Step 7
% Note: The below needs editing!!!

\renewcommand{\update}{
$
  \begin{array}{l}          % do not delete this line 
    \mbox{$C := A_1(A_1^T)^T + C$} \\ % replace \mbox{...} with update line but leave 
  \end{array}               % do not delete this line 
$
}




\begin{figure}
	\begin{center}
		\FlaWorksheet
	\end{center}
	\caption{Blocked Algorithm for Loop Invariant \eqref{eq:inv1}}
	\label{fig:blk_inv1}
\end{figure}


% \begin{figure}
% \includegraphics[width=\textwidth]{FLAMEC/Plot_unb_GFLOPS.pdf}
% \caption{Performance of the unblocked algorithm for loop invariant~\ref{eq:inv1} of TRMM}.
% \label{fig:unbAlg1}
% \end{figure}



\end{document}
